%==============================================================================
\chapter{Literature Survey}
\label{ch:literature}
%==============================================================================

This chapter surveys the six bodies of work that TutorAI draws on:
intelligent tutoring systems and AI in education; the cognitive theory of
learning from words and pictures; large language models as content
generators; agentic orchestration and self-correction of LLMs; automatic
diagram and vector-graphics generation; and speech synthesis together with
offline-first mobile architecture. Each section closes with the specific
lesson the surveyed work contributes to this thesis, and the chapter closes
with a summary table and the identified research gap.

%------------------------------------------------------------------------------
\section{Intelligent Tutoring Systems and AI in Education}
%------------------------------------------------------------------------------

The ambition of software that teaches is old and well studied. Bloom's
influential ``2 sigma'' finding --- that students taught one-to-one by a human
tutor performed about two standard deviations better than students in
conventional classrooms --- framed the field's goal: find scalable methods of
instruction as effective as individual tutoring~\cite{bloom1984}. Intelligent
Tutoring Systems (ITS) pursued this goal by modelling the domain, the student,
and pedagogy explicitly. Anderson \emph{et al.}'s cognitive tutors, grounded
in the ACT-R theory of cognition, demonstrated that model-tracing tutors could
achieve large learning gains in algebra and programming, and distilled a set
of pragmatic design lessons from a decade of classroom
deployment~\cite{anderson1995cognitive}. VanLehn's later meta-analysis
compared human tutoring, step-based ITS, and answer-based systems, and found
--- contrary to folklore --- that mature step-based tutoring systems were
nearly as effective as human tutors (effect sizes of roughly 0.76 versus
0.79)~\cite{vanlehn2011}. The catch has always been cost: classical ITS
require painstaking, per-domain knowledge engineering, which is why they
exist only for a handful of well-resourced subjects.

The arrival of large language models changed the economics of educational
content. Kasneci \emph{et al.} survey the opportunities and challenges of
LLMs such as ChatGPT in education: on the one hand personalized explanations,
instant feedback, and content generation at negligible marginal cost; on the
other hand hallucination, bias, over-reliance, and the difficulty of
verifying generated material~\cite{kasneci2023chatgpt}. This tension ---
generation is now cheap but \emph{trustworthy} generation is not --- is the
starting point of the present work.

\emph{Lesson for TutorAI.} TutorAI does not attempt a full ITS: it has no
student model and no adaptive curriculum. It automates one narrow, valuable
tutoring behaviour --- pointing at a drawing while explaining it --- for an
open set of topics, at near-zero marginal cost per repeated topic. Where a
classical ITS invests in hand-built domain knowledge, TutorAI invests in
\emph{verification} of machine-generated content.

%------------------------------------------------------------------------------
\section{Learning from Words and Pictures}
%------------------------------------------------------------------------------

Why insist on a diagram whose parts light up in time with speech, rather than
a paragraph of text? The answer comes from cognitive psychology. Paivio's
dual-coding theory holds that verbal and visual information are processed in
separate channels and that learning is stronger when both are
engaged~\cite{paivio1991}. Sweller's cognitive load theory adds that working
memory is sharply limited, so instruction should avoid forcing the learner to
hold and cross-reference disconnected representations~\cite{sweller1988}.

Mayer's cognitive theory of multimedia learning operationalizes these ideas
into evidence-backed design principles~\cite{mayer2014multimedia}. Three of
them read almost as a specification of TutorAI's player. The
\emph{multimedia principle}: people learn better from words and pictures than
from words alone. The \emph{modality principle}: spoken narration
accompanying a graphic beats on-screen text, because it spreads processing
across the auditory and visual channels. The \emph{temporal contiguity} and
\emph{signaling} principles: narration and the corresponding part of the
graphic should be presented at the same time, and the learner's attention
should be explicitly cued to the relevant part. Mayer and Moreno's survey of
ways to reduce cognitive load in multimedia learning names synchronized
narration-plus-highlighting as one of the most effective
techniques~\cite{mayermoreno2003}.

\emph{Lesson for TutorAI.} The product form --- segment-level narration,
each segment simultaneously highlighting exactly the diagram parts it
describes, with captions available --- is a direct implementation of the
multimedia, modality, temporal-contiguity, and signaling principles. This
grounding also explains why a \emph{broken} highlight (a narration segment
pointing at a non-existent part) is not a cosmetic bug but a destruction of
the pedagogical mechanism, and hence why the thesis treats reference
consistency as a hard correctness property.

%------------------------------------------------------------------------------
\section{Large Language Models for Content Generation}
%------------------------------------------------------------------------------

The transformer architecture~\cite{vaswani2017attention} and the subsequent
scaling of autoregressive language models established that a single
pre-trained model can perform diverse generation tasks from natural-language
instructions alone. Brown \emph{et al.} showed with GPT-3 that few-shot
prompting suffices for many tasks without task-specific
fine-tuning~\cite{brown2020gpt3}. The Gemini family extended this to natively
multimodal models spanning text, code, images, and audio, with ``Flash''-class
variants engineered for low latency and cost~\cite{gemini2023} --- the class
of model TutorAI uses for planning, drawing, narration, and repair. Beyond
prose, Chen \emph{et al.}'s Codex work demonstrated that LLMs trained on code
generate syntactically constrained, machine-interpretable artifacts with
useful reliability~\cite{chen2021codex}; generating SVG markup, as TutorAI
does, is essentially structured code generation. Liu \emph{et al.}
systematize prompting methods --- instruction design, output schemas,
temperature control --- as the practical interface for steering such
models~\cite{liu2023promptsurvey}.

Two properties of LLM generation matter for this thesis. First,
\emph{structured output}: modern APIs can constrain generation to a JSON
schema, which TutorAI uses for plans, narration segments, and repair
responses, validated by Pydantic on receipt. Second, \emph{hallucination}: Ji
\emph{et al.}'s survey documents that language models produce fluent content
untethered from their inputs across virtually all generation
tasks~\cite{ji2023hallucination}. Schema-constrained output fixes the
\emph{shape} of a response but not its \emph{referential truth}: narration can
be perfectly valid JSON and still reference the identifier
\texttt{mitochondria} when the drawing contains no such element.

For completeness, raster text-to-image models such as Stable
Diffusion~\cite{rombach2022stable} can also illustrate topics, but their
output is pixels: there are no addressable parts, so nothing can be
selectively highlighted, and no deterministic validator can check whether the
image matches the narration. This rules them out for TutorAI's purpose and
motivates vector output.

\emph{Lesson for TutorAI.} Use an LLM for what it is good at (planning,
drawing well-known concepts, fluent narration, structured JSON), but treat
referential consistency as a property to be \emph{checked deterministically},
never assumed.

%------------------------------------------------------------------------------
\section{LLM Agents, Orchestration, and Self-Correction}
%------------------------------------------------------------------------------

A second line of work addresses reliability by wrapping the model in a
\emph{process}. Chain-of-thought prompting showed that eliciting intermediate
reasoning steps improves accuracy on multi-step problems~\cite{wei2022cot}.
ReAct interleaved reasoning with tool-using actions, letting a model gather
facts and act on feedback~\cite{yao2023react}. Most directly relevant are the
self-correction methods: Self-Refine has the model critique and then improve
its own output iteratively~\cite{madaan2023selfrefine}; Reflexion converts
environment feedback into verbal self-reflections that improve subsequent
attempts~\cite{shinn2023reflexion}; and Chen \emph{et al.} teach models to
self-debug generated code against execution feedback~\cite{chen2023selfdebug}.
A consistent finding across these papers is that critique grounded in
\emph{concrete, externally produced error signals} (failing tests, parser
errors) outperforms free-form self-assessment.

Orchestration frameworks make such processes buildable. LangChain provides
the component model for prompts, models, and tools~\cite{langchain}, and
LangGraph structures an application as an explicit state graph --- typed
state flowing through nodes, with conditional edges for loops --- rather than
an open-ended autonomous loop~\cite{langgraph}. The graph form matters for
production systems: it makes control flow inspectable, each node separately
testable, and loop budgets explicit.

\emph{Lesson for TutorAI.} TutorAI adopts decomposition
(plan~$\to$~draw~$\to$~narrate), a bounded critique-and-refine loop for the
drawing, and a bounded repair loop driven by a validator's exact error list
--- the ``self-debugging with real feedback'' pattern. It deliberately
rejects open-ended agency: the graph is small, static, and every loop has a
configured budget, so worst-case cost and latency are capped. The novel
element relative to the surveyed work is that the final arbiter is not the
model's own judgement but a deterministic, non-LLM gate that enforces a
cross-artifact invariant before anything is cached.

%------------------------------------------------------------------------------
\section{Automatic Diagram and Vector-Graphics Generation}
%------------------------------------------------------------------------------

TutorAI needs drawings whose parts are \emph{addressable}. SVG, a W3C
standard, represents images as an XML document tree in which any element can
carry an \texttt{id} attribute and be styled or animated
individually~\cite{svgspec} --- precisely the property synchronized
highlighting requires, and one no raster format offers.

Prior work on generating vector graphics spans learned latent models and
direct code generation. DeepSVG learns a hierarchical generative network over
SVG path commands, targeting interpolation and animation of icon-scale
graphics~\cite{carlier2020deepsvg}. StarVector treats SVG as code and
generates it with a multimodal LLM from images or text, showing that
modern code models internalize the SVG grammar
well~\cite{rodriguez2023starvector}. General code-generation results
\cite{chen2021codex} support the same conclusion. However, none of these
systems constrain the \emph{semantics of identifiers}: they produce plausible
markup, not markup whose element ids form a contract with a separately
generated narration. Template- and grammar-based diagram tools (of which
Mermaid is a popular example) guarantee well-formedness by construction, but
only within a fixed diagram vocabulary (flowcharts, sequence diagrams), which
is too narrow for arbitrary educational topics like photosynthesis or the
water cycle.

\emph{Lesson for TutorAI.} Have the LLM draw free-form SVG (breadth), impose
an \emph{id contract} through prompting (every meaningful element gets a
semantic id; narration may reference only existing ids), and enforce the
contract with an XML parser rather than trusting the model (reliability).
Quality is then raised --- not guaranteed --- by the critique-and-refine
loop.

%------------------------------------------------------------------------------
\section{Speech Synthesis and Narrated Content}
%------------------------------------------------------------------------------

Neural TTS made synthesized narration acceptable for learning content.
WaveNet demonstrated autoregressive generation of raw audio waveforms with
naturalness far beyond concatenative systems~\cite{oord2016wavenet}, and
Tacotron~2 established the now-standard text~$\to$~mel-spectrogram~$\to$
vocoder pipeline achieving near-human mean opinion scores~\cite{shen2018tacotron2}.
Today, hosted TTS services (including Gemini's speech models) expose this
capability behind an API that returns audio for text.

The synchronization question --- how to keep visuals aligned with speech ---
is usually answered with word-level timestamps or forced alignment, which
requires either streaming infrastructure or per-word timing metadata, both
fragile on an offline mobile target. The alternative TutorAI adopts is
coarser but deterministic: synthesize one clip per narration segment,
\emph{measure} its exact duration from the PCM byte count and sample rate,
and drive highlighting from clip boundaries. Synchronization then requires no
timing data beyond what the device player already knows (which clip is
playing), and it survives offline replay trivially.

\emph{Lesson for TutorAI.} Segment-level synchronization trades pointer
granularity for determinism and offline robustness --- an explicit,
defensible design decision grounded in the target environment.

%------------------------------------------------------------------------------
\section{Offline-First Mobile Applications and Software Architecture}
%------------------------------------------------------------------------------

Offline-first design treats the network as an enhancement rather than a
prerequisite: applications keep a complete local copy of the data they need
and synchronize opportunistically~\cite{offlinefirst}. For a learning tool
aimed partly at users with intermittent connectivity, this is a requirement,
not a preference. On Android, Google's architecture guidance recommends a
layered structure --- UI, domain, and data layers with unidirectional data
flow and reactive state~\cite{androidguide} --- which TutorAI follows using
Jetpack Compose, ViewModels exposing immutable state via \texttt{StateFlow},
Room for structured storage, and a private file store for binary assets.

At the level of software structure, Martin's Clean Architecture argues for
dependency inversion: business rules at the centre, frameworks and I/O at the
edges, boundaries crossed through interfaces~\cite{martin2017clean}. The
classical design-patterns catalogue~\cite{gamma1994} supplies the vocabulary
TutorAI uses at its seams --- Strategy, Abstract Factory, Adapter, Builder,
Facade, Repository --- and the Twelve-Factor App methodology supplies the
operational discipline of strict configuration-through-environment
\cite{twelvefactor}, which is what lets the same codebase run on mock
providers in tests and real Gemini providers in production, and lets hosted
model identifiers be pinned in configuration rather than code.

\emph{Lesson for TutorAI.} None of these practices is novel individually;
the survey point is that an AI-centred system still succeeds or fails on
ordinary architectural discipline. The provider adapters, hermetic tests, and
configuration-first design described in Chapters~\ref{ch:methodology}
and~\ref{ch:design} are direct applications of this literature.

%------------------------------------------------------------------------------
\section{Summary and Research Gap}
%------------------------------------------------------------------------------

Table~\ref{tab:litsummary} condenses the survey.

\begin{table}[htbp]
\caption{Summary of the literature surveyed and its role in this thesis}
\label{tab:litsummary}
\centering
\small
\renewcommand{\arraystretch}{1.3}
\begin{tabular}{@{}p{0.26\textwidth}p{0.30\textwidth}p{0.34\textwidth}@{}}
\toprule
\textbf{Area (key works)} & \textbf{Established result} & \textbf{Use in TutorAI}\\
\midrule
ITS / AI in education
\cite{bloom1984,anderson1995cognitive,vanlehn2011,kasneci2023chatgpt}
& Tutoring software works but classical ITS need costly per-domain
engineering; LLMs make content cheap but unverified.
& Automate one tutoring behaviour (point-and-explain) for open topics;
invest in verification instead of hand-built domain knowledge.\\
Multimedia learning
\cite{paivio1991,sweller1988,mayer2014multimedia,mayermoreno2003}
& Synchronized narration + signaled graphics measurably improve learning.
& Defines the product form; makes reference consistency a hard
correctness property.\\
LLMs for generation
\cite{vaswani2017attention,brown2020gpt3,gemini2023,chen2021codex,liu2023promptsurvey,ji2023hallucination}
& LLMs generate prose, code, and structured data well, but hallucinate;
schemas fix shape, not referential truth.
& Gemini via adapters; structured-output prompts; motivates the
deterministic gate.\\
Agents / self-correction
\cite{wei2022cot,yao2023react,madaan2023selfrefine,shinn2023reflexion,chen2023selfdebug,langchain,langgraph}
& Decomposition and critique on concrete error signals improve
reliability; graphs make agent control flow explicit.
& LangGraph pipeline; bounded critique-and-refine; repair loop fed by
validator errors.\\
Vector-graphics generation
\cite{svgspec,carlier2020deepsvg,rodriguez2023starvector,rombach2022stable}
& LLMs can write SVG; raster models offer no addressable parts; grammar
tools are narrow.
& Free-form LLM-drawn SVG under a prompted id contract, enforced by an
XML parser.\\
TTS \cite{oord2016wavenet,shen2018tacotron2}
& Neural TTS is natural enough for narration; word-level sync needs
fragile timing data.
& Per-segment clips with byte-measured durations; segment-level sync.\\
Mobile / architecture
\cite{offlinefirst,androidguide,martin2017clean,gamma1994,twelvefactor}
& Offline-first, layered architecture, patterns, and 12-factor config are
proven practice.
& Clean Architecture + MVVM client; provider adapters; hermetic key-free
tests; container + CI/CD.\\
\bottomrule
\end{tabular}
\end{table}

\textbf{Research gap.} Each ingredient above exists in isolation: LLMs that
draw SVG, agent frameworks with self-correction, multimedia-learning theory
prescribing synchronized narration, and offline-first mobile practice. What
the surveyed literature does not provide is a system that \emph{combines}
them under a hard cross-artifact consistency guarantee --- narration that is
machine-checked, before caching, to reference only diagram elements that
exist --- delivered end to end on an offline-first mobile client with
deterministic audio--visual synchronization. Published self-correction work
validates outputs against tests or environments for a \emph{single} artifact;
TutorAI's gate validates the \emph{joint} consistency of two independently
generated artifacts (drawing and narration) and refuses service rather than
degrade. Closing that gap --- and demonstrating it with a verified, tested,
deployable implementation --- is the contribution of this thesis.
